\FloatBarrier
\section{本章小结}
\label{sec:region-integrals-summary}

区域积分把密度累计成总量。Jordan区域是能用有限个坐标盒从内外两侧任意逼近的有界集合，等价地，它的边界能被总体积任意小的有限个盒覆盖。边界因而不影响有界函数的累计：在区域闭包上连续的函数可积，分块连续的有界函数也可积；有界本身则不够，例如有理点的指示函数不可积。积分继承有限和的线性、保序、绝对值估计与区域可加性，体积、质量、平均值和质心都是区域积分。

累次积分把一次多维累计化为几次一元积分。被积函数连续、区域由连续的边界函数描述时，按截面分组的Riemann和收敛到累次积分，两种次序给出同一个值；换序时要重新求出区域在另一坐标轴上的投影和截面，必要时改变分块。区域无界或被积函数无界时，有界情形的结论不能直接使用：要先截断再取极限，而不同的截断方式和不同的积分次序可以给出不同的结果。

换元时，可微映射在一点附近近似线性，Jacobian行列式的绝对值给出局部体积的缩放，换元公式把这些倍数累计起来。公式要求两个开集之间的$C^1$双射及$C^1$逆，绝对值消去的是定向的符号，消不去覆盖次数。极坐标、柱坐标和球坐标的权重分别为$r$、$r$与$\rho^2\sin\varphi$；它们在原点、轴线和角度接缝处不满足条件，须先截去这些部分，再用被积函数的有界性控制遗漏的贡献。

参数积分在紧Jordan区域上依赖统一控制：被积函数在区域与参数集的乘积上连续时，积分随参数连续；参数偏导数也联合连续时，可以在积分号下求导，并由连续偏导数推出全微分。只有逐点的极限不够。积分限随参数移动时，导数除了积分号下的项，还要加上两端新增或移除的细条贡献。下一章把累计的范围从区域换成曲线。
